\documentclass[11pt]{article}
\usepackage{iftex}
\ifPDFTeX
  \usepackage[T1]{fontenc}
  \usepackage{lmodern}
\fi
\usepackage{amsmath,amssymb,amsthm}
\usepackage[margin=1in]{geometry}
\usepackage{booktabs,longtable,array}
\usepackage[numbers,sort&compress]{natbib}
\usepackage{microtype}
\usepackage{xurl}
\usepackage{hyperref}
\hypersetup{colorlinks=true,linkcolor=blue,citecolor=blue,urlcolor=blue,
  pdftitle={Exceptional sets for sums and differences of two primes along polynomial sequences},
  pdfauthor={Jizhou Guo}}
\Urlmuskip=0mu plus 1mu
\setlength{\emergencystretch}{2em}
\newtheorem{theorem}{Theorem}[section]
\newtheorem{lemma}[theorem]{Lemma}
\newtheorem{proposition}[theorem]{Proposition}
\newtheorem{corollary}[theorem]{Corollary}
\newtheorem{inputthm}[theorem]{Input}
\theoremstyle{definition}
\newtheorem{definition}[theorem]{Definition}
\newtheorem{hypothesis}[theorem]{Hypothesis}
\theoremstyle{remark}
\newtheorem*{remark}{Remark}
\numberwithin{equation}{section}
\newcommand{\e}{\mathrm e}
\newcommand{\eps}{\varepsilon}
\newcommand{\Z}{\mathbb Z}
\newcommand{\N}{\mathbb N}
\newcommand{\R}{\mathbb R}
\newcommand{\C}{\mathbb C}
\newcommand{\T}{\mathbb T}
\newcommand{\M}{\mathfrak M}
\newcommand{\m}{\mathfrak m}
\newcommand{\A}{\mathcal A}
\newcommand{\ind}{\mathbf 1}
\newcommand{\abs}[1]{\left\lvert#1\right\rvert}
\newcommand{\norm}[1]{\left\lVert#1\right\rVert}
\DeclareMathOperator{\Rea}{Re}
\DeclareMathOperator{\cond}{cond}
\title{Exceptional sets for sums and differences of two primes\protect\\
along polynomial sequences}
\author{Jizhou Guo\\
  \href{mailto:mitsuha2021b@gmail.com}{\texttt{mitsuha2021b@gmail.com}}}
\date{29 October 2026}

\begin{document}
\maketitle
\begin{abstract}
For a fixed integer polynomial $Q$ of degree $d=2$ or $3$ with positive
leading coefficient, we study the parameters $1\le m\le N$ for which
$Q(m)$ is positive and even but is not a sum of two primes.
Assuming the fourth-moment and major-arc theorems of Li--Liu under a
common fixed Dirichlet zero-free boundary $\Theta<11/12$, we obtain
$O_{Q,\eps}(N^{d/4+\eps})$ exceptions. The same bound holds for failure
of a prime-difference representation with both primes at most $C_QN^d$,
where $C_Q$ is explicit. Pintz's major-arc formula and the exponential-sum
estimate of Maynard--Pandey--Radziwi\l\l, under a common boundary
$\Theta<1$, give the alternative bound $O_{Q,\eps}(N^{7d/24+\eps})$.
For quadratic polynomials mapping an entire integer interval into
$[X/2,X]$, the corresponding bounds are $O_\eps(X^{1/4+\eps})$ and
$O_\eps(X^{7/24+\eps})$, uniformly in the coefficients and the interval.
OpenAI's seven-eighths theorem supplies the zero-free hypothesis in
both routes. The argument is the classical thin-set exceptional-set
mechanism, with the exponential sum taken over the exceptions themselves;
for $Q(t)=2t^2$ the alternative route is the argument of Perelli--Pintz
with this restriction. We give the elementary multiplicity arguments,
including the cubic zero branches.
\end{abstract}
\noindent\textit{2020 Mathematics Subject Classification.}
Primary 11P32; Secondary 11L07, 11M26.

\section{Statements and context}\label{sec:introduction}

All logarithms are natural, $e(t)=\exp(2\pi it)$, and primes may repeat.
Parameters $m\le N$ always mean positive integers $1\le m\le N$.
Let $Q\in\Z[t]$ have degree $d\in\{2,3\}$ and positive leading
coefficient. For $N,Z\ge2$ define
\begin{align}
 E_+(N;Q)
 &=\#\{1\le m\le N: Q(m)>0,\ 2\mid Q(m),\ Q(m)\ne p_1+p_2\},
       \label{eq:Eplus}\\
 E_-(N;Q;Z)
 &=\#\{1\le m\le N: Q(m)>0,\ 2\mid Q(m),\notag\\[-2mm]
 &\hspace{40mm} Q(m)\ne p_2-p_1
                         \text{ for primes }p_1,p_2\le Z\}.
       \label{eq:Eminus}
\end{align}
The coefficients of a fixed polynomial are independent of $N$.
Set
\begin{equation}\label{eq:CQ}
 A_Q=\sum_{j=0}^d\abs{[t^j]Q},\qquad C_Q=2(1+A_Q).
\end{equation}
We use the notation $\mathrm{ZF}(\Theta)$ for the following common
zero-free hypothesis: every nontrivial zero of every primitive
Dirichlet $L$-function, including the character of conductor one,
has real part at most the same fixed number $\Theta$.
The boundary is uniform in conductor and height. A principal pole at
$s=1$ is allowed. We state all analytic inputs in Section~\ref{sec:inputs}.

\begin{theorem}[Conditional on the Li--Liu inputs]\label{thm:fixed}
Assume $\mathrm{ZF}(\Theta)$ for a fixed $1/2\le\Theta<11/12$,
and use Li--Liu's Theorem~1 and Proposition~4.1 as stated in
Inputs~\ref{input:LL4} and~\ref{input:LLmajor}.
For every fixed $Q$ as above and every $\eps>0$,
\begin{equation}\label{eq:main}
 E_+(N;Q)+E_-(N;Q;C_QN^d)
       \ll_{Q,\Theta,\eps} N^{d/4+\eps}.
\end{equation}
Thus the parameter exponents are $1/2$ for quadratics and $3/4$ for
cubics. No irreducibility hypothesis or additional local condition
on $Q$ is required.
\end{theorem}

\begin{theorem}[Conditional Pintz--MPR alternative]\label{thm:fallback}
Assume $\mathrm{ZF}(\Theta)$ for a fixed $1/2\le\Theta<1$,
and use Pintz's Theorem~1, (2.10)--(2.11), and
Maynard--Pandey--Radziwi\l\l's Theorem~1.1 as stated in
Inputs~\ref{input:Pintz} and~\ref{input:MPR}.
For every fixed $Q$ as above and every $\eps>0$,
\begin{equation}\label{eq:fallback}
 E_+(N;Q)+E_-(N;Q;C_QN^d)
       \ll_{Q,\Theta,\eps} N^{7d/24+\eps}.
\end{equation}
The parameter exponents are $7/12$ and $7/8$, respectively.
\end{theorem}

Here and below constants may also depend on the fixed weight and the
fixed constants of the cited analytic inputs. The alternative theorem
uses a different prime sum and a different choice of arcs from
Theorem~\ref{thm:fixed}.

For a finite interval $I$ of consecutive integers and a quadratic
$Q\in\Z[t]$, put
\begin{align*}
 U_+(I;Q)&=\#\{m\in I:2\mid Q(m),\ Q(m)\ne p_1+p_2\},\\
 U_-(I;Q;Z)&=\#\{m\in I:2\mid Q(m),\ Q(m)\ne p_2-p_1
                                      \text{ with }p_1,p_2\le Z\}.
\end{align*}
The range condition in the next theorem makes all these values positive.
Perelli \cite[p.~478]{Perelli} leaves estimates uniform in the coefficients
as a separate question.

\begin{theorem}[Conditional uniform quadratic bound]\label{thm:uniform}
Let $X\ge2$, let $I$ be a finite interval of consecutive integers, and
let $Q\in\Z[t]$ have degree exactly two. Suppose that
\begin{equation}\label{eq:wholeinterval}
                 Q(m)\in[X/2,X]\qquad(m\in I).
\end{equation}
Under the inputs of Theorem~\ref{thm:fixed},
\begin{equation}\label{eq:uniformLL}
 U_+(I;Q)+U_-(I;Q;2X)\ll_{\Theta,\eps}X^{1/4+\eps}.
\end{equation}
Under the inputs of Theorem~\ref{thm:fallback},
\begin{equation}\label{eq:uniformfallback}
 U_+(I;Q)+U_-(I;Q;2X)\ll_{\Theta,\eps}X^{7/24+\eps}.
\end{equation}
The constants are independent of the coefficients of $Q$, the location
of $I$, and its length. Either sign separately is also bounded by $|I|$.
The leading coefficient may have either sign.
\end{theorem}

\begin{corollary}[Specialization of the zero-free input]\label{cor:openai}
OpenAI's Theorem~1.1, stated in Input~\ref{input:openai}, gives
$\mathrm{ZF}(7/8)$. Hence the conclusions of
Theorems~\ref{thm:fixed}--\ref{thm:uniform} hold with $\Theta=7/8$,
relative to the remaining inputs of each theorem
(Inputs~\ref{input:LL4}--\ref{input:LLmajor}, respectively
Inputs~\ref{input:Pintz}--\ref{input:MPR}). In particular this applies
to $Q(t)=2t^2$ and $Q(t)=2t^3$ and to their bounded prime-difference versions.
\end{corollary}

\subsection{The classical thin-set mechanism}

If $\A$ is the set of distinct exceptional polynomial values and
$E=|\A|$, the exponential sum
$K(\alpha)=\sum_{a\in\A}e(a\alpha)$ retains the exceptional set itself.
A bound $r_\A(s)\le D$ for ordered pair-sum multiplicities gives
\[
 \int_\T|K(\alpha)|^4\,d\alpha\le DE^2.
\]
The major coefficients are of size at least $X$, whereas the full
coefficients vanish at the exceptional values. H\"older's inequality
therefore transfers either a fourth moment or a pointwise bound for the
prime sum to a bound for $E$. Br\"udern--Kawada--Wooley
\cite[Introduction]{BKW} describe working with an exponential sum over
the exceptions. In the proof of their Theorem~1 \cite[\S2, (5)]{BKW},
as in Perelli--Pintz \cite[(3.6)--(3.10)]{PP}, this is the sign-weighted
sum over all parameters, whose moments are bounded by the full polynomial
mean value. Here the sum runs over the exceptional values only, so that
its fourth moment is at most $DE^2$.
The elementary part of the present argument supplies $D\ll X^\eta$
for fixed quadratics and cubics, and uniformly for the quadratic
families in \eqref{eq:wholeinterval}. Applying the two analytic
interfaces gives the bounds in Table~\ref{tab:exponents}.

\begin{table}[ht]
\centering
\begin{tabular}{@{}lccc@{}}
\toprule
Analytic route & Physical block & Fixed degree $2$ & Fixed degree $3$\\
\midrule
Li--Liu & $X^{1/4+\eps}$ & $N^{1/2+\eps}$ & $N^{3/4+\eps}$\\
Pintz--MPR & $X^{7/24+\eps}$ & $N^{7/12+\eps}$ & $N^{7/8+\eps}$\\
\bottomrule
\end{tabular}
\caption{Bounds for the number of exceptional parameters, relative to
the inputs in the corresponding theorem. A physical block means
$Q(m)\in[X/2,X]$; for fixed $Q$ its largest scale is $O_Q(N^d)$.
The physical bounds are uniform in the quadratic setting of
Theorem~\ref{thm:uniform}.}\label{tab:exponents}
\end{table}

The difference problem uses the Fourier coefficients of $|S|^2$ in
place of $S^2$. The fourth moment of $K$ still counts pair \emph{sums}
of exceptional values. This observation allows the same polynomial
multiplicity lemmas to treat both signs.

\subsection{Related work}\label{sec:related}

Perelli \cite[Theorem~1 and Corollary~1]{Perelli} treats a fixed
polynomial $F\in\Z[t]$ of degree $k$ with positive leading coefficient.
Writing his physical scale as $Y$, for
$Y^{1/(3k)+\eps}\le H\le Y^{1/k}$ his Corollary~1 gives, among the
parameters $Y^{1/k}\le m\le Y^{1/k}+H$ with even $F(m)$,
$O_{A,\eps,F}(H(\log Y)^{-A})$ Goldbach exceptions, for every $A>0$.
His Theorem~1 is the underlying mean-square estimate for the weighted
representation function. Br\"udern--Kawada--Wooley
\cite[Theorem~1]{BKW} prove that, for a fixed positive-leading
degree-$k$ polynomial $\Phi$, the number of $m\le N$ for which
$2\Phi(m)$ is not a sum of two primes is
$O_\Phi(N^{1-c/k})$, where $c>0$ is absolute. Their introduction
describes the use of an exponential sum over the exceptions; in the
proof of Theorem~1 this is a sign-weighted sum over all parameters
\cite[\S2, (5)]{BKW}. A polynomial
whose even values occur only on some residue classes is accommodated
by splitting $m$ modulo two: whenever $Q(r)$ is even,
$Q(2t+r)/2\in\Z[t]$.

For doubled squares write
\[
 E_2(N)=\#\{1\le m\le N:2m^2\ne p_1+p_2\}.
\]
Pintz \cite[Theorem~5]{Pintz2021}, in a theorem attributed there to
Perelli and Pintz, states, unconditionally, $E_2(N)\ll_\eps N^{4/5+\eps}$.
Perelli--Pintz \cite[Theorem~1 and (1.4)]{PP} state
$E_2(N)\ll_\eps N^{19/24+\eps}$, using the fixed Dirichlet zero-free
half-plane together with Pintz's major arcs and MPR's exponential-sum
estimate. Here $N$ is the parameter cutoff for $2m^2$. In physical
size $2m^2\le X$, their stated bound becomes $X^{19/48+\eps}$.
For $Q(t)=2t^2$ the proof of Theorem~\ref{thm:fallback} is the argument
of \cite[(3.1)--(3.10)]{PP}, with the same inputs and the same H\"older
inequality, except that the sign-weighted sum over all $m\le N$ is
replaced by the sum over the exceptional values. Its fourth moment is
then $O(E^2N^{\eps})$ instead of $O(N^{2+\eps})$, which gives
$N^{7/12+\eps}$ in place of $N^{19/24+\eps}$.
Theorem~\ref{thm:fixed} gives $N^{1/2+\eps}$ by using the Li--Liu
fourth moment in place of the pointwise bound.

Pintz--Alth\"ofer \cite[Theorem~1 and Corollary~2]{PA} prove
$E(X)\ll_\eps X^{7/12+\eps}$ for the ordinary Goldbach exceptional
set under any common fixed Dirichlet zero-free boundary below one,
and insert OpenAI's seven-eighths theorem. Their Sections~2.1--2.2
align Pintz's major arcs with the MPR bound. Li--Liu
\cite[Theorems~1--2 and Proposition~4.1]{LL} obtain
$E_+(X),E_-(X)\ll_\Theta X^{1/2}(\log X)^{100}$ under
$\mathrm{ZF}(\Theta)$ for $\Theta<11/12$; their new analytic estimate
is the minor-arc fourth moment. We use that estimate together with
their major coefficients, preserving its weight and arc parameter.

Later work lowers the zero-free boundary below $7/8$ (for instance
\cite{Gao,Argonaut,LiuRefinement,XuSang,LiuHecke}). The QRH tracker
records the Dirichlet boundary $0.8739558$ \cite{Sussman}.
Within the two routes of this paper such improvements leave every
exponent in Table~\ref{tab:exponents} unchanged: in the Li--Liu parameter
selection $\Theta$ occurs through $\max(\Theta,7/8)$, and in the
alternative route the term $X^{19/24}$ in the MPR estimate determines
the power $7/24$.

Section~\ref{sec:inputs} records the inputs and their compatibility.
Section~\ref{sec:transfer} proves the finite transfer, and
Section~\ref{sec:multiplicities} supplies the algebraic multiplicity
bounds. Sections~\ref{sec:LLproof}--\ref{sec:fallbackproof} prove the
main results. The appendix describes the scope of the ancillary checks.

\section{Analytic inputs}\label{sec:inputs}

\begin{hypothesis}[Common Dirichlet zero-free boundary]\label{hyp:ZF}
For one fixed $1/2\le\Theta<1$, every nontrivial zero $\rho$ of every
primitive Dirichlet $L$-function satisfies $\Rea\rho\le\Theta$.
This includes $\zeta$, the $L$-function of conductor one, and applies
at all conductors and all heights. Denote this assertion by
$\mathrm{ZF}(\Theta)$.
\end{hypothesis}

\begin{inputthm}[OpenAI, Theorem~1.1, Dirichlet specialization]\label{input:openai}
Every Dirichlet $L$-function, including $\zeta$, has no zero in
$\Rea s>7/8$. A pole at $s=1$ for a principal character is allowed.
This is the Dirichlet assertion of \cite[Theorem~1.1]{OpenAIQ78}.
\end{inputthm}

References to the OpenAI manuscript and its Lean scope document
\cite{OpenAIQ78,OpenAIScope} are pinned to the repository commit
\begin{center}\small
\nolinkurl{adc7f1241b42e322a6451854ab7e4b4c146bf78a}.
\end{center}
The full Theorem~1.1 also treats finite-order Hecke $L$-functions over
$\mathbb Q(\sqrt{-3})$. The Dirichlet assertion alone supplies
Hypothesis~\ref{hyp:ZF} at $\Theta=7/8$. The Dirichlet comparator
linked by the family~003 scope document uses the declaration name
\begin{center}\small
\nolinkurl{OAI.DirichletCharacter.LFunction_ne_zero_of_seven_eighths_lt_re}.
\end{center}
We use the written theorem with this scope. The ancillary computations
below verify finite algebra and examples.

\subsection{The Li--Liu weight, arcs, and theorems}

Fix a real function $w\in C_c^\infty((0,1))$ such that $0\le w\le1$
and $w=1$ on $[1/8,7/8]$. Choose it in the logarithmic Gevrey class
of order two specified in \cite[\S2]{LL}: for fixed constants
$C_0,C_1>0$ and all integers $j\ge0$,
\begin{equation}\label{eq:Gevrey}
 \sup_{v\in\R}\abs{\frac{d^j}{dv^j}w(e^v)}
             \le C_0C_1^j(j!)^2.
\end{equation}
The weight is fixed once and for all. Put
\begin{equation}\label{eq:prime-sum}
 S_X(\alpha)=\sum_p(\log p)w(p/X)e(p\alpha).
\end{equation}
On $\T=\R/\Z$ use Lebesgue measure of total mass one. Write
$\norm{t}_\T$ for the distance from $t$ to the nearest integer, and define
\begin{equation}\label{eq:arcs}
 \M_X(P)=\bigcup_{q\le P}\ \bigcup_{\substack{a\bmod q\\(a,q)=1}}
 \left\{\alpha:\norm{\alpha-a/q}_{\T}\le\frac{P}{qX}\right\},
 \qquad \m_X(P)=\T\setminus\M_X(P).
\end{equation}
For a fixed exponent $\kappa<1/2$ and $P=X^\kappa$, the arcs are
disjoint for sufficiently large $X$.

\begin{inputthm}[Li--Liu, Theorem~1]\label{input:LL4}
Assume $\mathrm{ZF}(\Theta)$ for a fixed $1/2\le\Theta<11/12$.
There is a fixed $\kappa=\kappa(\Theta)\in(4/9,9/20)$ such that
\begin{equation}\label{eq:LL4}
 \int_{\m_X(X^\kappa)}|S_X(\alpha)|^4\,d\alpha
                   \ll_{w,\Theta}X^{5/2}(\log X)^{100}.
\end{equation}
The weight and arcs are those in \eqref{eq:Gevrey}--\eqref{eq:arcs}.
\end{inputthm}

For a nonzero even integer $r$ define
\begin{equation}\label{eq:singularseries}
 \mathfrak S(r)=2\prod_{p>2}\left(1-\frac1{(p-1)^2}\right)
                  \prod_{\substack{p\mid r\\p>2}}\frac{p-1}{p-2}.
\end{equation}
The infinite product converges to a positive number, since
$\sum_{p>2}(p-1)^{-2}<\infty$, and each factor in the finite product
is at least one. Thus $\mathfrak S(r)\gg1$ uniformly in nonzero even $r$.
Define the finite weight convolutions
\begin{equation}\label{eq:Iw}
 I_w^+(n)=\sum_{k\in\Z}w(k/X)w((n-k)/X),\quad
 I_w^-(h)=\sum_{k\in\Z}w(k/X)w((k+h)/X).
\end{equation}

\begin{inputthm}[Li--Liu, Proposition~4.1]\label{input:LLmajor}
Assume $\mathrm{ZF}(\Theta)$ for a fixed $1/2\le\Theta<1$.
For every fixed $0<\kappa<9/20$ there is $c_1>0$ such that, with
$P=X^\kappa$,
\begin{align}
 \int_{\M_X(P)}S_X(\alpha)^2e(-n\alpha)\,d\alpha
   &=\mathfrak S(n)I_w^+(n)+O_{w,\Theta,\kappa}(X^{1-c_1}),
               \label{eq:LLmajorplus}\\
 \int_{\M_X(P)}|S_X(\alpha)|^2e(-h\alpha)\,d\alpha
   &=\mathfrak S(h)I_w^-(h)+O_{w,\Theta,\kappa}(X^{1-c_1}),
               \label{eq:LLmajorminus}
\end{align}
uniformly for even $n\in[X/2,X]$ and even $2\le h\le X/2$,
respectively. The weight is the same as in Input~\ref{input:LL4}.
\end{inputthm}

The flat portion of $w$ makes the main terms uniformly positive.
For $n\in[X/2,X]$, the intersection
$[X/8,7X/8]\cap[n-7X/8,n-X/8]$ has length at least $X/4$.
For $2\le h\le X/2$, the intersection
$[X/8,7X/8]\cap[X/8-h,7X/8-h]$ has length
$3X/4-h\ge X/4$. Counting integers in these intervals gives
$I_w^+(n),I_w^-(h)\ge X/4-O(1)$. Consequently, for some fixed $c>0$
and all sufficiently large $X$,
\begin{align}
 \Rea\int_{\M_X(P)}S_X^2 e(-n\alpha)\,d\alpha&\ge cX,
                    \label{eq:LLpositiveplus}\\
 \Rea\int_{\M_X(P)}|S_X|^2 e(-h\alpha)\,d\alpha&\ge cX
                    \label{eq:LLpositiveminus}
\end{align}
on the stated ranges. In particular Input~\ref{input:LLmajor} applies
to the same $\kappa$ supplied by Input~\ref{input:LL4}.

The proof of Li--Liu's Theorem~1 chooses $\kappa$ as in
\cite[Proposition~9.4]{LL}, which uses
\begin{equation}\label{eq:kappa-choice}
 \max\left\{\frac{1351}{3006},\frac{5-24\delta}{10-40\delta}\right\}
       <\kappa<\frac9{20},\qquad
 \delta=1-\max(\Theta,7/8)>\frac1{12}.
\end{equation}
At $\Theta=7/8$, taking $\kappa=899/2000$ gives
\begin{equation}\label{eq:kappa-exact}
 \frac{(5/2)(1/2-\kappa)}{1+10(1/2-\kappa)}=\frac{101}{1204}<\frac18,
 \qquad
 \frac92\left(\frac12-\kappa\right)=\frac{909}{4000}<\frac{38}{167}.
\end{equation}
The last margin is $197/668000$. These rational inequalities are
checked in \path{anc/exact_checks.py}. Theorem~1 of Li--Liu is used
with its existential choice of $\kappa$; \eqref{eq:kappa-choice}
records the narrower choice from its proof.

\subsection{Pintz's major-arc formula and MPR's pointwise estimate}

Fix a sufficiently small $\nu>0$, with $\nu<1/72$, and set
\begin{equation}\label{eq:P0}
 P_0=X^{4/9-\nu},\quad X_1=X^{1-\rho},\quad
 T_X(\alpha)=\sum_{X_1<p\le X}(\log p)e(p\alpha).
\end{equation}
Here $\rho>0$ is a fixed sufficiently small admissible cutoff exponent
chosen as in Pintz's theorem after $\nu$. We keep this cutoff exponent
separate from positive exponents in error terms.

\begin{inputthm}[Pintz, Theorem~1, (2.10)--(2.11), zero-free specialization]
\label{input:Pintz}
Use the numbering of \cite[arXiv v1]{PintzMajor}. Assume
$\mathrm{ZF}(\Theta)$ for a fixed $1/2\le\Theta<1$.
With \eqref{eq:P0}, for each fixed sufficiently large $H$ and all
sufficiently large $X$, one can choose the same
$P\in[P_0X^{-\nu},P_0]$ for the two formulas
\begin{align}
 \int_{\M_X(P)}T_X(\alpha)^2e(-n\alpha)\,d\alpha
 &=\mathfrak S(n)n+
 O_\nu(\mathfrak S(n)Xe^{-c_\nu H})+O_\nu(X^{1-c_2}),
                       \label{eq:Pintzplus}\\
 \int_{\M_X(P)}|T_X(\alpha)|^2e(-h\alpha)\,d\alpha
 &=\mathfrak S(h)(X-X_1-h+O(1))\notag\\
 &\quad+O_\nu(\mathfrak S(h)Xe^{-c_\nu H})+O_\nu(X^{1-c_2}),
                       \label{eq:Pintzminus}
\end{align}
for even $n\in[X/2,X]$ and even $h\in[X/4,X/2]$, respectively.
Here $c_\nu,c_2>0$ are fixed, and the implicit constants may depend
on the fixed parameters of the formula. The arcs are \eqref{eq:arcs}
at this same $P$.
\end{inputthm}

To specify this specialization, Pintz's set of generalized exceptional
singularities consists of the principal pole and nontrivial zeros with
conductor at most $P$, real part at least $1-H/\log X$, and ordinate
of absolute value at most $\sqrt X$. First fix $H$; then
$1-H/\log X>\Theta$ for large $X$, so only the pole remains.
Its term in (2.10) is $\mathfrak S(n)n$, since
$\Gamma(1)^2/\Gamma(2)=1$. In (2.11) its integer convolution is
\[
 \#\{(k,\ell)\in(X_1,X]^2\cap\Z^2:k-\ell=h\}
                  =X-X_1-h+O(1).
\]
The positive power $c_2$ may be decreased to absorb fixed logarithmic
factors or the cutoff error. Choosing $H$ large and then $X$ large
gives a common lower bound $cX$ for the real parts of both major
coefficients. This is a pointwise lower bound for every target in the
displayed ranges. The sum specialization also appears in
\cite[\S2.1]{PA}.

\begin{inputthm}[Maynard--Pandey--Radziwi\l\l, Theorem~1.1]\label{input:MPR}
Let $Y\ge2$, $q\in\N$, $(a,q)=1$, $q\le\sqrt Y$, and
$|\alpha-a/q|\le1/(q\sqrt Y)$. Put
$B_Y=\max\{q,qY|\alpha-a/q|\}$. For every $\eta>0$,
\begin{equation}\label{eq:MPR}
 \abs{\sum_{n<Y}\Lambda(n)e(n\alpha)}
       \ll_\eta Y^\eta\left(\frac{Y}{\sqrt{B_Y}}+Y^{19/24}\right).
\end{equation}
Here $\Lambda(p^j)=\log p$ for primes $p$ and integers $j\ge1$,
and $\Lambda(n)=0$ otherwise. This is the arbitrary-small-power
form of \cite[Theorem~1.1]{MPR}.
\end{inputthm}

Changing $n<Y$ to $n\le Y$ costs at most $\log(2Y)$.
Theorem~\ref{thm:fallback} uses the external estimate \eqref{eq:MPR}.

\section{A finite exceptional-set transfer}\label{sec:transfer}

For a finite set $\A$ of distinct integers write
\begin{equation}\label{eq:multiplicity}
 r_\A(v)=\#\{(a,b)\in\A^2:a+b=v\}.
\end{equation}
These pairs are ordered. For $K(\alpha)=\sum_{a\in\A}e(a\alpha)$,
finite orthogonality gives
\begin{equation}\label{eq:energy}
 \int_\T|K(\alpha)|^4\,d\alpha
       =\sum_{v\in\Z}r_\A(v)^2
       \le \left(\sup_v r_\A(v)\right)|\A|^2.
\end{equation}

\begin{lemma}[Finite transfer]\label{lem:transfer}
Let $S$ be a trigonometric polynomial on $\T$, and let $\M,\m$ be
complementary measurable subsets of $\T$. Let $\A$ be a finite set
of distinct integers, $E=|\A|$, and $c,X>0$. Choose either
$F=S^2$ or $F=|S|^2$. Suppose that, for every $a\in\A$,
\begin{equation}\label{eq:Fourierconditions}
 \int_\T F(\alpha)e(-a\alpha)\,d\alpha=0,\qquad
 \Rea\int_\M F(\alpha)e(-a\alpha)\,d\alpha\ge cX.
\end{equation}
If $r_\A(v)\le D$ for every integer $v$, put
$B_j=\int_\m|S|^j$ for $j=2,4$, and $M=\sup_\m|S|$, with
$M=0$ when $\m$ is empty. Then
\begin{equation}\label{eq:transfer}
 E\le c^{-2}X^{-2}D^{1/2}
              \min\{B_2B_4^{1/2},\ M B_2^{3/2}\}.
\end{equation}
\end{lemma}

\begin{proof}
If $E=0$, the conclusion holds directly. Suppose $E>0$, so $D\ge1$,
and put $K(\alpha)=\sum_{a\in\A}e(a\alpha)$. Summing the real parts
of the major coefficients in \eqref{eq:Fourierconditions} and using
the vanishing full coefficients gives
\begin{equation}\label{eq:basictransfer}
 cXE\le\abs{\int_\m F(\alpha)\overline{K(\alpha)}\,d\alpha}
              \le\int_\m|S(\alpha)|^2|K(\alpha)|\,d\alpha.
\end{equation}
Both choices of $F$ have modulus $|S|^2$. H\"older's inequality with
exponents $2,4,4$, followed by \eqref{eq:energy}, yields
\[
 cXE\le B_2^{1/2}B_4^{1/4}
                  \left(\int_\T|K|^4\right)^{1/4}
        \le B_2^{1/2}B_4^{1/4}D^{1/4}E^{1/2}.
\]
Alternatively, $|S|^2\le M^{1/2}|S|^{3/2}$ on $\m$. H\"older
with exponents $4/3,4$ gives
\[
 cXE\le M^{1/2}B_2^{3/4}
                  \left(\int_\T|K|^4\right)^{1/4}
        \le M^{1/2}B_2^{3/4}D^{1/4}E^{1/2}.
\]
Squaring either inequality and cancelling $c^2X^2E$ proves its part
of \eqref{eq:transfer}. If $\m$ is empty, the hypotheses force
$E=0$, which is covered by the first case.
\end{proof}

For either prime sum used here the second moment has the elementary
bound
\begin{equation}\label{eq:secondmoment}
 \int_\T\abs{\sum_{p\le X}(\log p)a_pe(p\alpha)}^2\,d\alpha
 =\sum_{p\le X}(\log p)^2|a_p|^2
 \le X(\log X)^2\qquad(|a_p|\le1).
\end{equation}
Thus the only prime-distribution estimates used in the transfer are
the major-arc lower bound and the chosen minor-arc input.

\begin{corollary}[The two physical powers]\label{cor:physical}
Suppose the conditions of Lemma~\ref{lem:transfer} hold, and for
every $\eta>0$ suppose $D\ll_\eta X^\eta$ and
$B_2\ll X(\log X)^2$, with the desired uniformity of constants.
If $B_4\ll_\eta X^{5/2+\eta}$, then
$E\ll_\eps X^{1/4+\eps}$. If instead
$M\ll_\eta X^{19/24+\eta}$, then
$E\ll_\eps X^{7/24+\eps}$.
\end{corollary}

\begin{proof}
The first expression on the right of \eqref{eq:transfer} is
\[
 \ll_\eta X^{-2+\eta/2}X(\log X)^2 X^{5/4+\eta/2}
        =X^{1/4+\eta}(\log X)^2.
\]
The second is
\[
 \ll_\eta X^{-2+\eta/2}X^{19/24+\eta}
                  X^{3/2}(\log X)^3
        =X^{7/24+3\eta/2}(\log X)^3.
\]
Choose $\eta$ sufficiently small relative to $\eps$, then absorb
the fixed logarithmic factors into the remaining small power.
\end{proof}

The set $\A$ in this application consists of the actual distinct
exceptions. Its cardinality $E$ enters the fourth moment as $E^2$,
preserving the arithmetic information used in the classical method.

\section{Polynomial pair-sum multiplicities}\label{sec:multiplicities}

\subsection{Divisors and quadratic polynomials}

\begin{lemma}[Elementary divisor bounds]\label{lem:divisors}
Let $\tau(n)$ be the number of positive divisors of the positive integer
$n$. For every $\eta>0$, $\tau(n)\ll_\eta n^\eta$.
For $V>0$, the number of ordered integer pairs satisfying
$x^2+y^2=V$ is at most $4\tau(V)$.
\end{lemma}

\begin{proof}
For all sufficiently large primes $p$ and every integer $a\ge1$,
$a+1\le2^a\le p^{\eta a}$. For each of the finitely many smaller
primes, $\sup_{a\ge0}(a+1)p^{-\eta a}<\infty$. Multiplying over the
prime factorization proves the divisor bound.

The Gaussian integers are Euclidean for the norm: after dividing in
$\C$, round both coordinates to obtain a remainder of norm at most
half the divisor's norm. The Euclidean algorithm gives a greatest
common divisor and B\'ezout's identity. An irreducible is therefore
prime, and induction on the norm gives unique factorization.
If a rational prime $p$ is reducible in $\Z[i]$, its factors have
norm $p$, since their integer norms multiply to $p^2$; thus
$p$ is associated to $\pi\bar\pi$. The two factors may be associates.
If $p$ remains irreducible, its exponent in an element of norm $V$
is determined by the exponent of $p$ in $V$, when such an element
exists. The same is true in the associate case. If $\pi,\bar\pi$
are nonassociate and $p^a\Vert V$, there are at most $a+1$ ways to
distribute their exponents. The four units give at most
$4\prod_{p^a\Vert V}(a+1)=4\tau(V)$ elements $x+iy$ of norm $V$.
\end{proof}

\begin{lemma}[Fixed quadratic]\label{lem:quadratic}
Fix $Q(t)=at^2+bt+c\in\Z[t]$ with $a\ne0$.
For $N\ge2$ and any $B\subseteq[1,N]\cap\Z$, the number of ordered
solutions $(u,v)\in B^2$ of $Q(u)+Q(v)=s$ is
$O_{Q,\eta}(N^\eta)$, uniformly in the integer $s$.
If instead $B$ is any set of integer parameters whose values lie
in $[0,X]$, $X\ge2$, that number is $O_{Q,\eta}(X^\eta)$,
independently of a parameter cutoff. Every value of $Q$ has at
most two parameter preimages.
\end{lemma}

\begin{proof}
The identity
\begin{equation}\label{eq:quadraticidentity}
       4aQ(t)=(2at+b)^2+4ac-b^2
\end{equation}
maps a solution injectively to $x=2au+b$, $y=2av+b$ satisfying
\begin{equation}\label{eq:quadraticnorm}
             x^2+y^2=4as+2b^2-8ac=:V.
\end{equation}
For $V<0$ there is no solution; for $V=0$ there is at most one;
for $V>0$ Lemma~\ref{lem:divisors} bounds their number by $4\tau(V)$.
When an original solution exists, $|s|\ll_QN^2$ and hence
$|V|\ll_QN^2$. Apply the divisor bound with exponent $\eta/2$.
In the physical-size setting $0\le s\le2X$, so
$|V|\ll_QX$ and the same proof applies with a suitably small
divisor exponent. Finally, $Q(t)-z$ has at most two roots for any $z$.
\end{proof}

\subsection{Fixed cubics and their zero branches}

\begin{lemma}[Fixed cubic]\label{lem:cubic}
Fix $Q(t)=at^3+bt^2+ct+d_0\in\Z[t]$ with $a\ne0$.
For $N\ge2$ and $B\subseteq[1,N]\cap\Z$, the number of ordered
solutions $(u,v)\in B^2$ of $Q(u)+Q(v)=s$ is
$O_{Q,\eta}(N^\eta)$, uniformly in the integer $s$.
If $B$ is any set of \emph{positive} integer parameters with
$Q(B)\subseteq[0,X]$, $X\ge2$, the number is
$O_{Q,\eta}(X^\eta)$, independently of a parameter cutoff.
Every value has at most three parameter preimages.
\end{lemma}

\begin{proof}
Put
\begin{equation}\label{eq:cubicAB}
 A=9ac-3b^2,\qquad B_0=2b^3-9abc+27a^2d_0.
\end{equation}
Direct expansion gives
\begin{equation}\label{eq:cubicidentity}
       27a^2Q(t)=(3at+b)^3+A(3at+b)+B_0.
\end{equation}
For $x=3au+b$, $y=3av+b$, the pair-sum equation implies
\begin{equation}\label{eq:cubicfactor}
 V:=27a^2s-2B_0=(x+y)(x^2-xy+y^2+A).
\end{equation}
The affine maps to $x$ and $y$ are injective.

If $V\ne0$, let $\ell=x+y$. This is a signed divisor of $V$;
there are $2\tau(|V|)$ possibilities. For fixed $\ell$,
\begin{equation}\label{eq:cubicproduct}
                 xy=\frac{\ell^2+A-V/\ell}{3}.
\end{equation}
The sum and product determine at most two ordered root pairs.
Thus there are at most $4\tau(|V|)$ solutions. For parameters at
most $N$ and an existing solution, $|V|\ll_QN^3$. For selected values
in $[0,X]$, one has $s\in[0,2X]$ and $|V|\ll_QX$.
Lemma~\ref{lem:divisors}, with the exponent decreased as needed,
gives both asserted estimates when $V\ne0$.

If $V=0$, the first branch in \eqref{eq:cubicfactor} is $x+y=0$.
It forces
\begin{equation}\label{eq:cubiclinearbranch}
                         u+v=-\frac{2b}{3a}.
\end{equation}
If the right side is a positive integer $H$, there are at most $H-1$
ordered positive pairs; otherwise there are none. This is a constant
depending only on the fixed polynomial. The other branch is
$x^2-xy+y^2=-A$. Since
\begin{equation}\label{eq:cubicpositiveform}
       2(x^2-xy+y^2)=(x-y)^2+x^2+y^2,
\end{equation}
it lies in the fixed disk $x^2+y^2\le-2A$ and has $O_Q(1)$ integer
points. For $A>0$ it is empty. These branches cover $V=0$;
their intersection causes only overcounting. The restriction to
positive parameters and the fixed coefficients are both used.
The three-preimage bound follows from the degree.
\end{proof}

For a set of distinct selected values, choose one parameter preimage
per value. This injects ordered value pairs into parameter pairs.
Hence Lemmas~\ref{lem:quadratic} and~\ref{lem:cubic} bound
\eqref{eq:multiplicity} for any subset of the selected values.

\subsection{Coefficient-uniform quadratics}

\begin{lemma}[Three consecutive values]\label{lem:uniform}
Under \eqref{eq:wholeinterval}, the pair-sum multiplicity of $Q(I)$
is $O_\eta(X^\eta)$, with a constant independent of the coefficients
and of $I$. The same holds for every subset of $Q(I)$, and every
value has at most two parameter preimages.
\end{lemma}

\begin{proof}
If $|I|\le2$, there are at most four ordered parameter pairs.
Otherwise let $m_0$ be the left endpoint and write
$q(j)=Q(m_0+j)=Aj^2+Bj+C$. The integers $j=0,1,2$ belong to the
translated interval. Set $y_j=q(j)$ for these three integers.
Then
\begin{align}
 A&=(y_2-2y_1+y_0)/2,\quad B=y_1-y_0-A,\quad C=y_0,
                      \label{eq:coefficientformulas}\\
 |A|&\le X/2,\qquad |B|\le X,\qquad |C|\le X.
                      \label{eq:coefficientbounds}
\end{align}
The translated polynomial has integer coefficients and $A\ne0$.
If $q(j)+q(k)=s$ with parameters in this interval, then $s\in[X,2X]$.
Its quadratic norm parameter satisfies
\begin{equation}\label{eq:uniformnorm}
                     |4As+2B^2-8AC|\le10X^2.
\end{equation}
Use the cases $V<0$, $V=0$, $V>0$ from \eqref{eq:quadraticnorm}.
In the last case the divisor bound with exponent $\eta/2$ gives a
constant depending only on $\eta$. Passing to distinct value pairs,
and then to subsets, decreases the count. Each value has at most
two preimages by the degree.
\end{proof}

The interval condition supplies the three consecutive values.
The proof allows an interval crossing the vertex and a negative
leading coefficient. For example $q(j)=-j^2+2j+1$ on $j=0,1,2$
has values $1,2,1$ in $[1,2]$, including a repeated value.

\begin{proposition}[Moving cubic multiplicities]\label{prop:moving}
For integers $M\ge3$ let
\begin{equation}\label{eq:moving}
            Q_M(t)=2(2t-M-1)^3+6M^3,\qquad X=8M^3.
\end{equation}
The values $Q_M(1),\ldots,Q_M(M)$ are even, positive, distinct, and
belong to $[X/2,X]$. Their pair-sum multiplicity at $12M^3$ is at
least $M$. Consequently a coefficient-independent bound
$O_\eta(X^\eta)$ for these multiplicities fails for every
$0<\eta<1/3$.
\end{proposition}

\begin{proof}
For $1\le t\le M$, $z_t=2t-M-1$ lies in $[-(M-1),M-1]$. Thus
\[
 4M^3\le6M^3-2(M-1)^3\le Q_M(t)
     \le6M^3+2(M-1)^3\le8M^3.
\]
The cubic is strictly increasing, and the values are even.
Since $z_{M+1-t}=-z_t$, their sum is $12M^3$ for all $M$ ordered
choices of $t$. Finally $M/(8M^3)^\eta\to\infty$ when $\eta<1/3$.
\end{proof}

This construction exhibits the growing linear zero branch for moving
cubic coefficients. Its conclusion concerns the uniform multiplicity
extension of Lemma~\ref{lem:uniform}. The constant for a fixed cubic
in Lemma~\ref{lem:cubic} is allowed to depend on $Q$.

\section{The Li--Liu consequences}\label{sec:LLproof}

We make explicit the Fourier coefficients to which the transfer is
applied. By finite orthogonality,
\begin{align}
 \int_\T S_X(\alpha)^2e(-n\alpha)\,d\alpha
 &=\sum_{p_1+p_2=n}(\log p_1)(\log p_2)
                    w(p_1/X)w(p_2/X),\label{eq:fullplus}\\
 \int_\T|S_X(\alpha)|^2e(-h\alpha)\,d\alpha
 &=\sum_{p_2-p_1=h}(\log p_1)(\log p_2)
                    w(p_1/X)w(p_2/X).\label{eq:fullminus}
\end{align}
The second identity assigns the positive exponent to $p_2$.
In both sums the support of $w$ places the primes below $X$.
For a value having no prime-sum representation,
\eqref{eq:fullplus} is zero. For a value having no prime-difference
representation with both primes at most $X$, \eqref{eq:fullminus}
is zero. These implications supply the vanishing conditions in
Lemma~\ref{lem:transfer}.

\subsection{Physical blocks for sums}

Assume the inputs of Theorem~\ref{thm:fixed}. Fix $Y$ above the
analytic threshold and let $\A$ consist of the distinct even values
$Q(m)\in(Y/2,Y]$, $m\ge1$, that are not sums of two primes.
The physical forms of Lemmas~\ref{lem:quadratic} and~\ref{lem:cubic}
give
\begin{equation}\label{eq:Dfixed}
                           r_\A(v)\ll_{Q,\eta}Y^\eta.
\end{equation}
Use $S_Y$ with the common $\kappa$ of Inputs~\ref{input:LL4}
and~\ref{input:LLmajor}. Equations~\eqref{eq:LLpositiveplus},
\eqref{eq:secondmoment}, and~\eqref{eq:LL4} verify all the
conditions of Corollary~\ref{cor:physical}. Hence
\begin{equation}\label{eq:LLphysical}
                         |\A|\ll_{Q,\Theta,\eps}Y^{1/4+\eps}.
\end{equation}
There are at most $d$ parameters per value, so the same
bound holds for parameters on this physical block.

\begin{proof}[Proof of the sum assertion in Theorem~\ref{thm:fixed}]
For $1\le m\le N$, $|Q(m)|\le A_QN^d=:T$.
Use the disjoint physical blocks
\begin{equation}\label{eq:anchoredblocks}
                         (T2^{-j-1},T2^{-j}],\qquad j\ge0.
\end{equation}
For fixed $\gamma>0$, apply \eqref{eq:LLphysical} with exponent
$1/4+\gamma$ to the blocks above the analytic threshold. Their
sum is
\[
 \ll_{Q,\Theta,\gamma}\sum_{j\ge0}(T2^{-j})^{1/4+\gamma}
               \ll_{Q,\Theta,\gamma}T^{1/4+\gamma}.
\]
Values below the threshold form a fixed finite set of integers,
each with at most $d$ preimages, so they contribute $O_{Q,\Theta}(1)$.
Choose $\gamma>0$ so that $d\gamma\le\eps$ and substitute $T=A_QN^d$.
This gives the sum bound in \eqref{eq:main}.
\end{proof}

\subsection{One global scale for differences}

\begin{proof}[Proof of the difference assertion in Theorem~\ref{thm:fixed}]
Take the prime scale $Z=C_QN^d$. Every positive even value
$h=Q(m)$ with $m\le N$ lies in $[2,Z/2]$, by \eqref{eq:CQ}.
Let $\A$ consist of those distinct values which fail to have a
prime-difference representation with both primes at most $Z$.
Use $S_Z$, $F=|S_Z|^2$, and the same Li--Liu arcs.
Equation~\eqref{eq:fullminus} gives the vanishing full coefficient,
and \eqref{eq:LLpositiveminus} gives its major lower bound.
The physical multiplicity lemmas apply to selected values in
$[0,Z]$ and give $r_\A(v)\ll_{Q,\eta}Z^\eta$ for the entire set.
Corollary~\ref{cor:physical} yields
$|\A|\ll_{Q,\Theta,\gamma}Z^{1/4+\gamma}$.
Multiply by at most $d$ preimages per value and choose
$d\gamma\le\eps$. The bounded range of $Z$ below the analytic
threshold is absorbed by increasing the constant.
\end{proof}

The smooth sum delivers primes strictly below $Z$. The common
statement uses $\le Z$ to include the endpoint convention in the
alternative route. An integer having no unrestricted prime-difference
representation belongs to this bounded exceptional set as well.

\subsection{Uniform quadratic families}

\begin{proof}[Proof of \eqref{eq:uniformLL}]
Assume first that $X$ exceeds the analytic threshold. For sums take
the distinct exceptional even values in $Q(I)$, use $S_X$, and replace
\eqref{eq:Dfixed} by Lemma~\ref{lem:uniform}. Corollary~\ref{cor:physical}
gives $O_{\Theta,\eps}(X^{1/4+\eps})$ values, uniformly in the
coefficients. Each value has at most two parameter preimages.
For differences take prime scale $2X$. All target values lie in
$[(2X)/4,(2X)/2]$, within the range of Input~\ref{input:LLmajor}.
The same multiplicity estimate and the difference transfer give
the same count, and the Fourier support supplies primes below $2X$.

For bounded $X$ there are at most $\lfloor X\rfloor+1$ integer values
in $[X/2,X]$, each with at most two preimages. Their total is absorbed
by a constant depending only on the fixed analytic threshold and
$\eps$, uniformly in $Q$ and $I$. This includes empty intervals and
those of length at most two.
\end{proof}

\section{The Pintz--MPR alternative}\label{sec:fallbackproof}

The minor-arc estimate must hold for the same lower-truncated prime
sum $T_X$ and the same $P$ as the major coefficients in
Input~\ref{input:Pintz}.

\begin{lemma}[Both prefixes on the same minor arcs]\label{lem:compatible}
Fix the parameters of Input~\ref{input:Pintz}. For its chosen $P$,
\begin{equation}\label{eq:MPRminor}
               \sup_{\alpha\in\m_X(P)}|T_X(\alpha)|
                            \ll_\eta X^{19/24+\eta}
\end{equation}
for every $\eta>0$, assuming Input~\ref{input:MPR}.
\end{lemma}

\begin{proof}
We have
\begin{equation}\label{eq:Prange}
            X^{4/9-2\nu}\le P\le X^{4/9-\nu},
                       \qquad 4/9-2\nu>5/12.
\end{equation}
Fix $\alpha\notin\M_X(P)$. For each $Y\in\{X,X_1\}$ choose a
Dirichlet approximation $a/q$ with $q\le\sqrt Y$ and
$|\alpha-a/q|\le1/(q\sqrt Y)$, reducing the fraction if needed.
The approximation is taken modulo one in \eqref{eq:arcs}.
If $q>P$, then $B_Y\ge q>P\ge PY/X$. If $q\le P$, exclusion
from the same major arcs gives $|\alpha-a/q|>P/(qX)$, and again
$B_Y>PY/X$. Thus
\begin{equation}\label{eq:prefixineq}
 \frac{Y}{\sqrt{B_Y}}<\sqrt{\frac{XY}{P}}\le\frac{X}{\sqrt P}
                      \le X^{7/9+\nu}\le X^{19/24}.
\end{equation}
Apply \eqref{eq:MPR} to both prefixes, changing the strict endpoint
if needed, and subtract. Since $Y\le X$, both the $Y^{19/24}$ term
and the arbitrary small power are bounded at scale $X$.
The total weight of proper prime powers is
\begin{equation}\label{eq:primepowers}
          \sum_{\substack{p^j\le X\\j\ge2}}\log p
                         \ll X^{1/2}(\log X)^2.
\end{equation}
Indeed there are $O(\log X)$ possible exponents, at most $\sqrt X$
bases for each, and weight at most $\log X$ per term.
Removing these terms leaves $T_X$ and costs less than the bound in
\eqref{eq:MPRminor}.
\end{proof}

The small constant $\nu$ is chosen admissibly for the major-arc
formula before $X$ varies. As a rational illustration of the exponent
comparison, $\nu=1/144$ gives $P_0=X^{7/16}$,
$P\ge X^{31/72}$, and $X/\sqrt P\le X^{113/144}<X^{19/24}$.
The proof only requires the strict inequality in \eqref{eq:Prange}.

\begin{lemma}[Alternative physical block]\label{lem:fallbackphysical}
Under the inputs of Theorem~\ref{thm:fallback}, for a fixed quadratic
or cubic $Q$, the number of positive integer parameters with even
$Q(m)\in(Y/2,Y]$ unrepresented as a prime sum is
$O_{Q,\Theta,\eps}(Y^{7/24+\eps})$.
The number lacking a prime-difference representation with both primes
at most $2Y$ has the same bound. For a quadratic family satisfying
\eqref{eq:wholeinterval} at scale $Y$, both bounds have constants
independent of the coefficients and the parameter interval.
\end{lemma}

\begin{proof}
For sums use $T_Y$, its arcs from Input~\ref{input:Pintz}, and the
distinct exceptional values. The full Fourier coefficient is zero,
the major real part is at least $cY$, and
\eqref{eq:secondmoment} gives $B_2\le Y(\log Y)^2$.
Lemma~\ref{lem:compatible} gives $M\ll_\eta Y^{19/24+\eta}$.
The physical polynomial multiplicity lemmas and the second part of
Corollary~\ref{cor:physical} give the desired bound for values;
the degree bounds their parameter preimages.

For differences use prime scale $X=2Y$. Every target
$h\in(Y/2,Y]$ lies in $[X/4,X/2]$, exactly the range used in
\eqref{eq:Pintzminus}. Failure of a representation with both primes
at most $X$ makes the full coefficient of $|T_X|^2$ vanish.
The same transfer gives the asserted count and prime bound.
For uniform quadratics use Lemma~\ref{lem:uniform} instead of the
fixed-coefficient lemmas. Bounded scales are covered by the number
of integer values and the degree bound, as in the proof of
\eqref{eq:uniformLL}.
\end{proof}

\begin{proof}[Proof of Theorem~\ref{thm:fallback} and
\eqref{eq:uniformfallback}]
For fixed $Q$ sum Lemma~\ref{lem:fallbackphysical} over
\eqref{eq:anchoredblocks} with $T=A_QN^d$. The geometric sum is
$O_{Q,\Theta,\gamma}(T^{7/24+\gamma})$ for either sign.
For each difference block its delivered primes have size at most
$2T2^{-j}\le2A_QN^d<C_QN^d$. Every parameter failing the global
prime bound therefore belongs to a block failure just counted.
Values below the fixed analytic threshold give boundedly many
parameters. Choose $d\gamma\le\eps$ to obtain \eqref{eq:fallback}.
For a uniform quadratic family apply the physical lemma at $Y=X$,
giving primes at most $2X$ and \eqref{eq:uniformfallback}.
\end{proof}

\begin{proof}[Proof of Corollary~\ref{cor:openai}]
Input~\ref{input:openai} supplies Hypothesis~\ref{hyp:ZF} at
$\Theta=7/8$, below both required boundaries.
Apply the preceding theorems with this zero-free input and the
remaining inputs specified in each theorem.
\end{proof}

\appendix
\section{Ancillary checks}\label{app:checks}

The ancillary programs check the polynomial identities
\eqref{eq:quadraticidentity}, \eqref{eq:cubicidentity}, and
\eqref{eq:cubicpositiveform}; the rational parameter choice recorded
in \eqref{eq:kappa-exact}; and the exponent comparisons in the
finite transfer and in \eqref{eq:prefixineq}.
The program \path{anc/exact_checks.py} checks the identities and
the fixed and moving cubic examples. The program
\path{anc/interface_checks.py} checks the coefficient reconstruction,
the power comparisons, and the change from $19/24$ to $7/12$ when
the Perelli--Pintz sum is restricted to the exceptions.
The program \path{anc/verify_claims.py} additionally enumerates
pair-sum multiplicities of fixed quadratics and cubics on finite
parameter ranges and checks the three-value bounds on integer samples.

These are exact finite checks of the algebra and examples used here.
The deductions for general parameters are proved in
Sections~\ref{sec:transfer}--\ref{sec:fallbackproof}; the analytic
theorems are the inputs of Section~\ref{sec:inputs}.
The ancillary README gives the commands, dependencies, and scope of
each check.

\bibliographystyle{plainnat}
\bibliography{references}

\section*{Use of AI}
The proofs and the text of this paper were produced with AI systems
(GPT-6.1 Sol, GPT-6 Astra and Claude Opus~5.5) under the author's
direction, checked by independent reviews with AI models and by the
computations in the ancillary files. No human expert has reviewed
them. The author takes responsibility for the content.
\end{document}
